% --- Sprint backlog — même style visuel que le product backlog (chapitre 2) ---
% Contenu des user stories / tâches 1.x–6.x issu du tableau Word « Tableau 4.1 » ;
% tâches 1.7–1.10 ajoutées pour refléter le cycle plan côté API (soumission, validation, suppression, bulk).
\small
\PFETableSetup

\setlength{\PFEbacklogTabW}{\dimexpr\textwidth+2.8cm\relax}
\setlength{\LTcapwidth}{\PFEbacklogTabW}
\setlength{\LTleft}{\fill}
\setlength{\LTright}{\fill}

\setlength{\LTpre}{0pt}
\setlength{\LTpost}{0pt}

\PFETableArrayStretch
\setlength{\extrarowheight}{4pt}
\renewcommand{\arraystretch}{1.4}

\arrayrulecolor{black}

\begin{longtable}{|
>{\centering\arraybackslash}m{1.2cm}|
>{\raggedright\arraybackslash}m{5.2cm}|
>{\centering\arraybackslash}m{1.6cm}|
>{\raggedright\arraybackslash}m{5.5cm}|
>{\centering\arraybackslash}m{1.7cm}|}

\caption{Sprint backlog du sprint 2 — Plans \& activités}\label{tab:sprint-backlog-2}\\

\hline
\tableHeaderStyle
\tableHeaderCell{ID} & \tableHeaderCell{User story} & \tableHeaderCell{Tâche ID} & \tableHeaderCell{Tâche} & \tableHeaderCell{Priorité} \\
\hline
\endfirsthead

\hline
\multicolumn{5}{|c|}{\tablename\ \thetable{} — \textit{(suite)}} \\
\hline
\tableHeaderStyle
\tableHeaderCell{ID} & \tableHeaderCell{User story} & \tableHeaderCell{Tâche ID} & \tableHeaderCell{Tâche} & \tableHeaderCell{Priorité} \\
\hline
\endhead

\hline
\multicolumn{5}{|r|}{\textit{Suite page suivante\ldots}} \\
\hline
\endfoot

\hline
\endlastfoot

1 & En tant qu'utilisateur (Site ou Corporate), je veux créer un plan CSR afin de planifier les activités & 1.1 &
Implémenter le module de gestion des plans CSR & +++++ \\
\cline{3-5}
& & 1.2 &
Permettre à l'utilisateur site de créer un plan CSR pour les sites auxquels il est rattaché & +++++ \\
\cline{3-5}
& & 1.3 &
Permettre à l'utilisateur corporate de créer un plan CSR à l'échelle de tous les sites & +++++ \\
\cline{3-5}
& & 1.4 &
Définir les informations du plan (site, année, budget total, \ldots) & +++++ \\
\cline{3-5}
& & 1.5 &
Modifier un plan CSR avant validation & ++++ \\
\cline{3-5}
& & 1.6 &
Consulter les plans CSR par année / statut / site & +++ \\
\cline{3-5}
& & 1.7 &
Soumettre un plan CSR pour validation (passage au statut soumis) & +++++ \\
\cline{3-5}
& & 1.8 &
Permettre à l'utilisateur corporate de valider ou de rejeter un plan soumis avec motif & +++++ \\
\cline{3-5}
& & 1.9 &
Supprimer un plan CSR lorsque les règles métier le permettent & ++++ \\
\cline{3-5}
& & 1.10 &
Soumettre ou supprimer plusieurs plans en une seule opération (actions en masse) & ++++ \\
\cline{3-5}
& & 1.11 &
Rechercher un plan CSR (par mot-clé, site, année) & ++++ \\
\cline{3-5}
& & 1.12 &
Filtrer les plans CSR selon des critères (statut, site, année) & ++++ \\
\cline{3-5}
& & 1.13 &
Soumettre une demande de modification sur un plan verrouillé & ++++ \\
\hline
2 & En tant qu'utilisateur, je veux importer un plan CSR via un fichier Excel afin de gagner du temps & 2.1 &
Implémenter la fonctionnalité d'import Excel & +++++ \\
\cline{3-5}
& & 2.2 &
Lire et parser le fichier Excel avant l'importer & +++++ \\
\cline{3-5}
& & 2.3 &
Valider les données (site, année, mode de validation) pour l'utilisateur site & +++++ \\
\cline{3-5}
& & 2.4 &
Détecter les conflits (doublons, incohérences) & +++++ \\
\cline{3-5}
& & 2.5 &
Afficher un aperçu avant l'importer & ++++ \\
\cline{3-5}
& & 2.6 &
Supprimer les duplicatas avant la validation & ++++ \\
\cline{3-5}
& & 2.7 &
Ignorer les duplicatas et continuer avant la validation & ++++ \\
\hline
3 & En tant qu'utilisateur Site, je veux gérer les activités planifiées afin d'organiser les actions CSR & 3.1 &
Ajouter une méthode pour ajouter activité planifiée & +++++ \\
\cline{3-5}
& & 3.2 &
Ajouter les informations nécessaires de chaque activité planifiée(titre,numéro, catégorie, budget prévu, Impact de l'action prévue \ldots) & +++++ \\
\cline{3-5}
& & 3.3 &
Ajouter une méthode de  modification de l'activité planifiée avant validation & ++++ \\
\cline{3-5}
& & 3.4 &
Ajouter une méthode de suppression de l'activité planifiée avant validation & ++++ \\
\cline{3-5}
& & 3.5 &
Ajouter une méthode de consultation du détail d'une activité planifiée & ++++ \\
\cline{3-5}
& & 3.6 &
Ajouter une fonctionnalité de recherche des activités planifiées & ++++ \\
\cline{3-5}
& & 3.7 &
Ajouter une fonctionnalité de filtrage des activités planifiées & ++++ \\
\hline
4 & En tant qu'utilisateur Site, je veux saisir les activités réalisées afin de suivre l'exécution & 4.1 &
Implémenter le module des activités réalisées & +++++ \\
\cline{3-5}
& & 4.2 &
Ajouter une méthode pour saisir  les données de l'activité réalisée( Budget réalisé, participants, Date de réalisation, \ldots) & +++++ \\
\cline{3-5}
& & 4.3 &
Ajouter une méthode de  modification de l'activité réalisée & ++++ \\
\cline{3-5}
& & 4.4 &
Ajouter une méthode de suppression de l'activité réalisée & ++++ \\
\cline{3-5}
& & 4.5 &
Ajouter une méthode de consultation du détail d'une activité réalisée & ++++ \\
\cline{3-5}
& & 4.6 &
Ajouter une fonctionnalité de recherche des activités réalisées & ++++ \\
\cline{3-5}
& & 4.7 &
Ajouter une fonctionnalité de filtrage des activités réalisées & ++++ \\
\hline
5 & En tant qu'utilisateur Site, je veux ajouter des activités hors plan afin de tracer les actions imprévues & 5.1 &
Ajouter une méthode de création d'une activité hors plan & +++++ \\
\cline{3-5}
& & 5.2 &
Saisir toutes les informations de l'activité hors plan (titre, catégorie, date de réalisation, \ldots) & +++++ \\
\cline{3-5}
& & 5.3 &
Ajouter une méthode de  modification de l'activité hors plan & ++++ \\
\cline{3-5}
& & 5.4 &
Ajouter une méthode de suppression de l'activité hors plan & ++++ \\
\cline{3-5}
& & 5.5 &
Soumettre une activité hors plan pour validation & ++++ \\
\hline
6 & En tant qu'utilisateur, je veux visualiser les informations planifiées et réalisées afin d'avoir une vue complète & 6.1 &
Afficher les données planifiées( Budget prévu, impact prévu de l'action, nombre de participants prévus, \ldots) & +++++ \\
\cline{3-5}
& & 6.2 &
Afficher les données réalisées( Budget réalisé, impact de l'action , nombre réel de participants, \ldots) & +++++ \\
\hline

\end{longtable}

\normalsize
